\section{Предметная область}

В качестве предметной области выбран подбор параметров шрифтовой гарнитуры для бейджей, изготавливаемых методом офсетной печати. Бейдж имеет ограниченную площадь, но при этом должен быстро считываться пользователем: на нём обычно размещаются имя, логотип организации и, при необходимости, должность, подразделение или дополнительная служебная информация.

Основная задача подбора шрифта заключается в том, чтобы сохранить читаемость текста при заданных условиях просмотра и печати. На читаемость влияют размер символов, контраст текста и фона, насыщенность начертания, форма знаков, межбуквенный интервал и различимость тонких элементов при визуальном восприятии.

Основой для выбора критериев служит стандарт ISO~24509:2019, описывающий оценку минимального читаемого размера символов с учётом условий просмотра, в том числе дистанции чтения, яркости и контраста~\cite{iso24509}. Для задачи печати бейджей эти факторы были адаптированы к практическим признакам макета и материала. В системе учитываются следующие критерии:

\begin{enumerate}
    \item наличие дополнительного текста помимо имени и логотипа;
    \item наличие кириллицы в текстовом содержимом;
    \item необходимость чтения бейджа с расстояния более одного метра;
    \item использование тёмного фона;
    \item наличие ламинации;
    \item плотность бумаги выше 250 г/м\textsuperscript{2}.
\end{enumerate}

Наличие дополнительного текста повышает требования к разборчивости гарнитуры, поскольку на той же площади нужно разместить больше информации. Кириллица учитывается отдельно, так как кириллические символы сложнее считывать с увеличенной дистанции. Увеличенная дистанция чтения требует более выраженных форм символов, большей высоты строчных знаков и повышенного контраста. Тёмный фон и ламинация дополнительно усложняют восприятие текста: блики и снижение контраста делают тонкие штрихи менее надёжными. Высокая плотность бумаги, наоборот, позволяет использовать более аккуратные детали.

Множество возможных решений системы образуют рекомендации по следующим параметрам шрифта:

\begin{itemize}
    \item тип гарнитуры:
    \begin{itemize}
        \item гротеск;
        \item гуманистический гротеск;
        \item антиква;
    \end{itemize}

    \item насыщенность начертания:и
    \begin{itemize}
        \item Regular;
        \item SemiBold;
        \item Bold;
    \end{itemize}

    \item допустимость тонких элементов:
    \begin{itemize}
        \item допускаются;
        \item допускаются умеренно;
        \item следует избегать;
    \end{itemize}

    \item высота строчных знаков:
    \begin{itemize}
        \item стандартная;
        \item увеличенная;
    \end{itemize}

    \item межбуквенный интервал:
    \begin{itemize}
        \item стандартный;
        \item слегка увеличенный;
        \item увеличенный;
    \end{itemize}

    \item контраст текста и фона:
    \begin{itemize}
        \item достаточный;
        \item высокий.
    \end{itemize}
\end{itemize}
